\subsubsection{Détermination constructive et fonction du registre dodécaphasique}
\label{alpha-ampl-dependencias}

La définition précédente permet de préciser ce que signifie déterminer alpha à partir
de la structure. Le domaine de l'opération contient quatre registres de
provenance connue : les trois suites régionales et le registre entier
de phase. La normalisation incorpore en outre une condition de frontière
compatible avec le prolongement. Une fois ces données fixées, la division
euclidienne détermine chaque reste et son quotient. La proposition
\ref{prop:alpha-normalizacion} étend cette détermination au système
complet. L'indépendance à l'égard d'une valeur physique cible est donc
une propriété de l'ordre de construction : la valeur comparée à la
fin n'intervient dans la sélection d'aucun des arguments de cette
application.

Cette indépendance causale doit être distinguée de l'indépendance algébrique
entre nombres. La première s'établit en identifiant le domaine, les
opérations et la provenance des coefficients ; la seconde serait une
assertion relative à des relations polynomiales. Le résultat utilisé ici est le
premier. La construction détermine une coordonnée au moyen de registres
produits par APP--TRIT--TPK et permet de vérifier ultérieurement les
relations qu'elle satisfait. Cet ordre permet d'utiliser une coordonnée déjà
engendrée dans une opération ultérieure sans la transformer en paramètre externe.

Le registre $K$ possède une fonction plus précise que celle d'une correction
numérique ajoutée à une formule. Ses douze positions proviennent de la
transformation réversible du registre signé. L'équation
$U=\mathcal H_{12}K$ conserve la possibilité de revenir à ce registre,
tandis que les différences transversales et la somme totale permettent une autre
reconstruction. Dans l'évaluation périodique, l'ordre des positions
détermine le poids positionnel de chaque composante. Le même ensemble de
composantes, disposé dans un autre ordre, produit en général une autre coordonnée
rationnelle. L'origine de phase et l'orientation font donc partie de la
définition du nombre qui intervient dans \eqref{eq:alpha-identidad-completa}.

La normalisation d'alpha compose cette information avec les régions. Son
identité locale distingue deux opérations successives. On forme d'abord le
bilan orienté des blocs ; ce bilan est ensuite représenté par
un reste canonique et un quotient reliant des positions contiguës. La première
opération conserve le signe du bilan et permet l'extraction ultérieure
d'un support. La seconde permet de former les préfixes de la coordonnée
réelle. Les deux lectures restent liées parce que l'état conserve les
données antérieures et postérieures à la normalisation. Le support signé et le
développement positionnel sont ainsi des réalisations différentes d'une composition
arithmétique explicite.

L'identité télescopique exprime le contenu global de cette composition.
Les quotients intérieurs s'annulent par paires lors de l'évaluation d'un préfixe,
tandis que leurs extrémités demeurent dans l'égalité. Cette annulation est une
propriété de la somme pondérée ; elle n'élimine pas les quotients de l'état et
ne permet pas de les supprimer avant d'effectuer la somme. La frontière droite d'une
fenêtre reçoit la contribution de la continuation. Pour cette raison,
l'évaluation d'une fenêtre isolée et la restriction d'une section
prolongée sont des opérations distinctes. Le passage du bloc final $663$ à
$662$, démontré précédemment, fournit une réalisation explicite de cette
différence sans altérer les blocs déjà stabilisés à sa gauche.

En ce sens, la bidirectionnalité possède une formulation arithmétique
concrète. La génération fournit des prolongements compatibles ; la
normalisation transmet l'information entière de leur frontière depuis les
positions de moindre poids vers celles de poids supérieur. L'équation conserve la
contribution des deux directions au moyen de ses indices et de ses conditions
aux limites. La réversibilité dans la fibre fixée se vérifie en reconstruisant
$K$ à partir des autres registres et des quotients complets. Cette
inversion certifie la cohérence de l'opération d'origine ; la direction
généalogique reste celle qui produit d'abord le registre puis
la coordonnée d'alpha.

La seconde détermination possède un contenu complémentaire. L'objet
immédiat y est un bilan de lacunes, et le paramètre apparaît comme
sa racine dans un domaine isolant. La voie positionnelle fournit une
normalisation canonique des registres ; la voie analytique étudie l'annulation
d'un opérateur. L'exposé des deux exige de conserver précisément
cette différence de fonctions. Leurs relations avec l'état commun doivent
s'exprimer au moyen des applications de raffinement et de la comparaison de
leurs réalisations complètes, comme le développe la sous-section
correspondante. L'accord d'une approximation finie constitue un
contrôle localisé ; l'identification des constructions complètes
possède le domaine et la portée de son propre énoncé.

Enfin, cette organisation explique la position d'alpha dans l'article.
Son équation apparaît après la génération régionale et la construction
de $K$, parce qu'elle utilise les deux résultats. La section électronique utilise
ensuite cette coordonnée comme un facteur de sa composition. Le prolongement
exceptionnel reçoit les registres et leurs supports, où l'information
d'incidence est conservée. Ces deux continuations reprennent une même
provenance et développent des fonctions mathématiques distinctes : évaluation d'une
section centrale et construction d'une configuration de frontière.
Le développement de chacune permet de suivre la structure à travers
l'équation, au lieu de limiter la comparaison à sa valeur scalaire.
